\documentclass[12pt,twoside]{article}
\usepackage{fancyhdr,makeidx,times,multicol,geometry}
\usepackage{keyval,ifthen,mparhack,manyfoot,poemscol}
\newcounter{scounter}
\newcommand\step{\stepcounter{scounter}\Roman{scounter}}

\title{Postcard from Heaven}
\author{J.J. Gómez Cadenas}
\begin{document}
\maketitle

\poemtitle{\step} 
\begin{poem}
%\begin{stanza} 
%\step
%\end{stanza}
\begin{stanza} 
We spent the first night\\
in a disgusting dump, whose name,\\
{\em Mission Universitaire Française,}\\
did not advertise the common bedroom, smelling\\
of bodies crammed into hammocks.\\
But we were twenty,\\
and even that sleepless night in the barracks\\
we found romantic.
\end{stanza}
\begin{stanza} 
Then a week, longer \\
than the last forty years. Wandering\\
around the Pompidou\\
and its crowded square. I remember\\
music, tightrope walkers,\\ 
mimes motionless as statues,\\ 
and the endless hours by the fountain,\\
talking about ART,\\
with the capital letters of youth.
\end{stanza}
\begin{stanza}
Yes, and the fancy coloured tubes,\\
bright and fuzzy, like a Kandinsky,\\
crossing the facade of the museum, evoking\\
the illusion of a false factory, as if a dream.
\end{stanza}
\begin{stanza}
Amazing to realise that once\\
I too was happy in Paris.
\end{stanza}
\begin{stanza}
Not that the city has changed so much since then.\\
The Seine is still there,\\
along with the stalls of old booksellers\\
and the streets of the Jewish quarter.
\end{stanza} 
\begin{stanza}
And I still haunt the usual places. Orsay,\\
Rodin's house, the Pompidou,\\
the long avenues of the Latin Quarter.\\
Always something to do,\\
a good dinner to enjoy,\\
a good friend to visit.
\end{stanza} 
\begin{stanza}
%But I never returned with you, and anyway,\\
But you never return twice to the same city.
\end{stanza} 
\begin{stanza}
You are twenty only once, naive enough to believe\\
that the city was freshly built,\\
just for us. It didn’t cross my mind what a cliché it was,\\
kissing you on the banks of the river, under the moon. Then,\\
we rushed to the hotel in Montmartre,\\
with its air of brothel barely remodelled,\\
that enormous, pretentious, maroon carpet,\\
and the giant mirror that repeated\\
the entanglement of our lithe bodies.
\end{stanza}
\begin{stanza}
We sneaked into the subway. Stole postcards,\\
read Baudelaire and Cavafy to each other,\\
and the city indulged us, complacent.\\ 
It was the middle of May,\\
and we believed that spring\\
would last forever.
\end{stanza} 
\end{poem}

\poemtitle{\step} 
\begin{poem}
%\begin{stanza} 
%\step
%\end{stanza}
\begin{stanza} 
The first city\\
exists before you ever set foot in Paris.\\
It is made of the stuff of dreams,\\
its long avenues stretching in your mind,\\
its arrogant buildings,\\
evoked by words.\\
Words heard in a thousand conversations,\\
read in a million novels,\\
repeated by poets and singers,\\
turned into images by filmmakers.\\
The first city is imagined, and yet\\
you will always ---in vain---\\
try to return there. Like Ithaca,\\
Paris never was. It was conjured\\
by a dirty river, a burned cathedral, catacombs,\\
greed, lust, empty stones,\\
and longing.
\end{stanza}
\end{poem}

\poemtitle{\step} 
\begin{poem}
%\begin{stanza} 
%\step
%\end{stanza}
\begin{stanza} 
You did not visit the second city either,\\
the one that you roamed at twenty,\\
madly in love, without realizing\\
that you were inventing both love and Paris.
\end{stanza} 
\begin{stanza} 
Those long days,\\
walking through streets as new as your eyes,\\
those feverish nights,\\
almost touching Nirvana,\\
the cafes which made you feel an artist,\\
the occasional stolen sandwich, playing bohemian,\\
%beer consumed by litres\\
%---it had a sour taste and let your mouth dry and thick,\\
%you didn't like very much drinking,\\ 
%or smoking those Marijuana joins,\\
%but it came with the job and so you abided---,\\
%and above it all,\\
the moon shining above your never-ending kisses.
\end{stanza} 
\begin{stanza} 
All of that\\
was just the plot of a soap opera.\\
The city was crowded with people repeating the same lines,\\
and yet no one realised. All the other million couples,\\
going through the same words,\\ 
strolling the same parks,\\
kissing like a disciplined army by the river,\\
%fucking in identical overpriced ex-brothels,\\
all were invisible to you, under the blinding light\\
of your twenty years. You never travel to Paris at that age.\\
It could have been Vienna, or Rome, Tel Aviv or Prague.\\
You just travel to the city of your youth. 
\end{stanza}
\end{poem}

\poemtitle{\step} 
\begin{poem}
%\begin{stanza} 
%\step
%\end{stanza}
\begin{stanza} 
Then many cities follow,\\
each one piling upon the last,\\
ordered in time, like archaeological layers,\\
and at the same time overlapping,\\
as if they were made of mist.
\end{stanza} 
\begin{stanza} 
You return to Paris so often over the years,\\
quick visits, work trips, conferences,\\
each time renewing your covenant with the place.
\end{stanza}
\begin{stanza} 
The magic is gone, yet\\
a good feeling persists.\\
You become a friend of the same streets,\\
and buildings, and museums, and bookstores,\\
that once were your lovers.
\end{stanza}
\begin{stanza} 
As if meeting an old flame,\\
there is between you both the intimacy,\\
which only gives the memory of naked bodies\\
under clean sheets. Yet the passion has passed.\\
You see wrinkles and decay where once\\
you only saw desire.
\end{stanza}
\end{poem}

\poemtitle{\step} 
\begin{poem}
%\begin{stanza} 
%\step
%\end{stanza}
\begin{stanza} 
But then you return to the city with your children,\\
and all the angels of Heaven\\
rush to greet you.
\end{stanza} 
\begin{stanza} 
You walk Paris again. They are by your side.\\
Your girl is eleven, and she carries the map\\
to sail the hostile underground and ensure\\
you all get to Disney. Her mom has told her,\\
Daddy won't find his way. He is half blind\\
and can't navigate the dark waters\\
of the subway.
\end{stanza} 
\begin{stanza} 
So she is taking care of you, \\
holding your hand firmly. But your son's grip\\
is even stronger. He is seven. His job\\
is to be your bodyguard. His stern gaze\\
scares away trolls and goblins. 
\end{stanza}
\begin{stanza} 
Suddenly,\\
your eyes are open again,\\
your heart is pounding with a love\\
you can hardly survive. \\
Not that selfish invention of yesteryear,\\
but something so deep and burning\\
that the city becomes real\\
for the first time.   
\end{stanza} 
\end{poem}

\poemtitle{\step} 
\begin{poem}
%\begin{stanza} 
%\step
%\end{stanza}
\begin{stanza} 
Eventually you return to Paris,\\
as an old man. There is something in the air\\
that rejects you. The old covenant, good for so long,\\
is torn to pieces.
\end{stanza} 
\begin{stanza} 
What is it? It takes you a while to realise the reason.\\
You still move around,\\
watching everything and everybody,\\
absorbing land and people alike\\
with your searching eyes.
\end{stanza}
\begin{stanza} 
When you were younger,\\
that burning stare\\
was a weapon. Gazes were crossed, like swords,\\
but no blood was spilled. Girls\\
looked back, sometimes.\\
Sparks flew in the night.
\end{stanza}
\begin{stanza} 
Now, if you are not careful,\\
you will be taken for a loony,\\
or worse, a lecher. \\
You need to remind yourself\\
that old men should avoid\\ 
eye contact with the young,\\
smile benignly ---trying not to drool---,\\
and hope to be spared. 
\end{stanza} 
\end{poem}

\poemtitle{\step} 
\begin{poem}
%\begin{stanza} 
%\step
%\end{stanza}
\begin{stanza} 
But then, again,\\ 
you return to the city with your children.
\end{stanza}
\begin{stanza} 

Your daughter is twenty, still holds your hand,\\
talks non-stop about her latest love-boy.\\
He is coming to meet her.\\
They will walk the same streets,\\
and kiss in the same corners\\
that once smiled at you.\\
They are inventing, one more time, \\
the same old tale,\\
yet so new.\\
The same old lies,\\
yet so true. 
\end{stanza}
\begin{stanza} 
Your son is eighteen,\\
wants to be a scientist like Mom and Dad.\\
When he was small,\\
he wanted to work in the same place\\
as both of you. Now\\
he is moving to Paris, away from home,\\
to fulfil his dreams.  
\end{stanza} 
\begin{stanza} 
How strange, to feel so happy,\\
and so hopeless. So proud of him,\\
as a king watching his prince being knighted, yet,\\
a black hole eating your heart. 
\end{stanza} 
\begin{stanza} 
Turns out, \\
love was not all those will-o'-the-wisps,\\
vanishing in the night,\\
like the light of fireflies. 
\end{stanza}

\begin{stanza} 
Love was something else,\\
made of clay and ashes,\\
made of blood and tears,\\
made of guilt and forgiveness, \\
made of years. 
\end{stanza}

\begin{stanza} 
Turns out,\\
Paris was not Paradise,\\
but it can still be\\
a postcard from Heaven.
\end{stanza}
\end{poem}

\poemtitle{Paris and Alexandria}
\begin{poem}
\begin{stanza}
So must I do as Antony did?\\
Open the windows of my coming old age,\\
and hear the divine troupe,\\
with exquisite music and laughter,\\
pass under my window?
\end{stanza}
\begin{stanza}
Shall I say goodbye to Paris's hotels?\\
To the beautifully ugly flophouse of our first night,\\
and to the room full of carpets and mirrors, high in Montmartre,\\
to that rundown hotel, glorious in its decay,\\
where an old waiter asked us, each morning,\\
if we had enjoyed our croissant?
\end{stanza}
\begin{stanza}
Shall I say goodbye\\
to the packed places where the four of us slept,\\
each one identical to the others,\\
in those years when time did not seem to flow?
\end{stanza}
\begin{stanza}
Shall I say goodbye to my children's laughter,\\
to the music that played\\
in our car, even more packed\\
than our hotel rooms,\\
with Mom and Dad and Children and Bird,\\
oh yes, and Bacchus too,\\
the God who shared our rides and our walks\\
through the streets of Paris at night?
\end{stanza}
\begin{stanza}
Shall I say goodbye? Is that all?\\
Part with Paris,\\
as Antony parted with Alexandria?
\end{stanza}
\begin{stanza}
Shall I accept, bravely, without regret, the inevitable departure of the God?
\end{stanza}
\begin{stanza}
Or say no, refuse, cling like a shipwrecked man to the last floating wood of love,
\end{stanza}
\begin{stanza}
and rage, rage against the dying of the light?
\end{stanza}
\end{poem}

\poemtitle{The clochard}
\begin{poem}
\begin{stanza}
The clochard who sleeps in the doorway\\
of that building on Boulevard du Montparnasse\\
owns a dirty mattress covered by a dirty rag,\\
a pillow,\\
and four books.
\end{stanza}
\begin{stanza}
Often, he sits perched, like an old vulture,\\
on the subway grates,\\
barefoot, white-bearded, a bottle of vodka at hand,\\
warmed by the alcohol and the fumes from the underground.\\
His books keep him company.
\end{stanza}
\begin{stanza}
Sometimes he stares at nothingness,\\
as if searching for the meaning of life.\\
Sometimes he reads,\\
royal and languid as a lord in his library.
\end{stanza}
\begin{stanza}
Nearby, tourists\\
spend on a ticket to the modernist museum\\
what he needs to survive for a full month.
\end{stanza}
\end{poem}
\end{document}
